
\begin{center}
{\LARGE\bfseries Kontrolllösung zum Arbeitsblatt: Excel-Tabellen}\\[1mm]
{\large Klasse 7 -- Schulfest-Kiosk planen und auswerten}
\end{center}

\begin{boxK}
Diese Kontrolllösung dient dem Abgleich nach der eigenen Bearbeitung. Die kurze Excel-Hilfe ist auch hier enthalten, weil die Lernenden ihre Fehler mithilfe der Begriffe und Befehle nachvollziehen sollen.
\end{boxK}

\section*{Kurze Excel-Hilfe}
\begin{tabularx}{\textwidth}{|L{3.0cm}|Y|L{3.0cm}|}
\hline
\rowcolor{gray!20}\textbf{Begriff} & \textbf{Bedeutung} & \textbf{Beispiel}\\\hline
Arbeitsmappe & die ganze Excel-Datei & Schulfest-Kiosk.xlsx\\\hline
Tabellenblatt & eine einzelne Seite in der Datei & Messreihe, Schueler\_Vorlage\\\hline
Spalte & senkrechter Bereich & A, B, C\\\hline
Zeile & waagerechter Bereich & 1, 2, 3\\\hline
Zelle & Schnittpunkt aus Spalte und Zeile & C3\\\hline
Zellbezug & Verweis auf eine bestimmte Zelle & C3 bedeutet: Wert aus C3\\\hline
\end{tabularx}

\begin{boxHinweis}
\code{=} startet eine Formel. \code{*} bedeutet mal. \code{-} bedeutet minus. \code{G3:G8} bedeutet alle Zellen von G3 bis G8. \code{SUMME} addiert. \code{MITTELWERT} berechnet den Durchschnitt. \texttt{\$B\$13} bleibt beim Herunterziehen fest.
\end{boxHinweis}

\section*{Etappe 0: Orientierung}
\textbf{Ziel:} Mit der Tabelle soll herausgefunden werden, welche Produkte beim Schulfest den größten Gewinn bringen und welche Produkte man beim nächsten Mal besonders einplanen sollte.

\textbf{Kontrollfrage:} Wichtiger ist der größte Gewinn, nicht nur der höchste Verkaufspreis. Der Gewinn berücksichtigt auch Einkaufskosten und verkaufte Anzahl.

\section*{Etappe 1: Excel-Grundlagen}
\begin{tabularx}{\textwidth}{|L{3.1cm}|Y|L{3.5cm}|}
\hline
\rowcolor{gray!20}\textbf{Begriff} & \textbf{Bedeutung} & \textbf{Beispiel}\\\hline
Spalte & senkrechter Bereich & A, B, C\\\hline
Zeile & waagerechter Bereich & 1, 2, 3\\\hline
Zelle & Schnittpunkt aus Spalte und Zeile & C3\\\hline
Zellbezug & Verweis auf eine Zelle & C3\\\hline
Formel & Rechenbefehl für Excel & \code{=C3*D3}\\\hline
\end{tabularx}

Eine Formel beginnt mit \code{=}, damit Excel die Eingabe als Rechnung erkennt.

\textbf{Kontrollfragen:} a) C3, b) \code{*}, c) alle Zellen von G3 bis G8.

\section*{Etappe 2: Messwerte eintragen}
\begin{center}
\begin{tabular}{|c|c|c|c|}
\hline
 & \textbf{Spalte 1} & \textbf{Spalte 2} & \textbf{Spalte 3}\\\hline
\textbf{Zeile 1} & 8 & 10 & 12\\\hline
\textbf{Zeile 2} & 7 & 11 & 15\\\hline
\textbf{Zeile 3} & 9 & 13 & 14\\\hline
\end{tabular}
\end{center}
\textbf{Kontrollfragen:} a) 15, b) 13, c) Zeile 1, Spalte 3.

\section*{Etappe 3: Kiosk-Tabelle lesen}
\begin{tabularx}{\textwidth}{|L{4cm}|Y|}
\hline
\rowcolor{gray!20}\textbf{Größe} & \textbf{Erklärung}\\\hline
Einkaufspreis & So viel kostet ein Produkt im Einkauf pro Stück.\\\hline
Verkaufspreis & So viel zahlen Kundinnen und Kunden pro Stück.\\\hline
Verkaufte Anzahl & So oft wurde der Artikel verkauft.\\\hline
Umsatz & Verkaufspreis mal verkaufte Anzahl.\\\hline
Kosten & Einkaufspreis mal verkaufte Anzahl.\\\hline
Gewinn & Umsatz minus Kosten.\\\hline
\end{tabularx}
\textbf{Kontrollfragen:} a) Spalte C, b) Verkaufspreis und verkaufte Anzahl.

\section*{Etappe 4: Umsatz berechnen}
\begin{tabularx}{\textwidth}{|L{5cm}|Y|}
\hline
\rowcolor{gray!20}\textbf{Frage vor der Formel} & \textbf{Antwort}\\\hline
Welche Werte brauche ich? & Verkaufspreis und verkaufte Anzahl.\\\hline
In welchen Zellen stehen sie beim Wasser? & C3 und D3.\\\hline
Welche Rechnung passt? & Verkaufspreis mal Anzahl.\\\hline
Welche Formel kommt in E3? & \code{=C3*D3}\\\hline
\end{tabularx}

\begin{tabularx}{\textwidth}{|Y|L{3cm}|Y|L{3cm}|}
\hline
Wasser & 86,00\,\euro{} & Brezel & 108,00\,\euro{}\\\hline
Apfelschorle & 76,80\,\euro{} & Muffin & 81,00\,\euro{}\\\hline
Eistee & 63,80\,\euro{} & Obstbecher & 72,60\,\euro{}\\\hline
\end{tabularx}
\textbf{Kontrollfrage:} Wasser: 86,00\,\euro{}.

\section*{Etappe 5: Kosten und Gewinn}
\begin{tabularx}{\textwidth}{|L{4.1cm}|Y|}
\hline
\rowcolor{gray!20}\textbf{Formel} & \textbf{Erklärung}\\\hline
Kosten in F3: \code{=B3*D3} & Einkaufspreis mal verkaufte Anzahl.\\\hline
Gewinn in G3: \code{=E3-F3} & Umsatz minus Kosten.\\\hline
\end{tabularx}

\begin{tabularx}{\textwidth}{|Y|L{2.7cm}|L{2.7cm}|}
\hline
\rowcolor{gray!20}\textbf{Artikel} & \textbf{Kosten} & \textbf{Gewinn}\\\hline
Wasser & 30,10\,\euro{} & 55,90\,\euro{}\\\hline
Apfelschorle & 32,00\,\euro{} & 44,80\,\euro{}\\\hline
Eistee & 26,10\,\euro{} & 37,70\,\euro{}\\\hline
Brezel & 50,40\,\euro{} & 57,60\,\euro{}\\\hline
Muffin & 36,00\,\euro{} & 45,00\,\euro{}\\\hline
Obstbecher & 36,30\,\euro{} & 36,30\,\euro{}\\\hline
\end{tabularx}
\textbf{Kontrollfragen:} Gewinn Wasser: 55,90\,\euro{}. Gewinn Brezel: 57,60\,\euro{}.

\section*{Etappe 6: Summen und Mittelwert}
\begin{tabularx}{\textwidth}{|L{3.5cm}|L{4.2cm}|Y|}
\hline
\rowcolor{gray!20}\textbf{Zelle} & \textbf{Formel} & \textbf{Ergebnis}\\\hline
D10 & \code{=SUMME(D3:D8)} & 358\\\hline
E10 & \code{=SUMME(E3:E8)} & 488,20\,\euro{}\\\hline
F10 & \code{=SUMME(F3:F8)} & 210,90\,\euro{}\\\hline
G10 & \code{=SUMME(G3:G8)} & 277,30\,\euro{}\\\hline
B12 & \code{=MITTELWERT(G3:G8)} & ca. 46,22\,\euro{}\\\hline
\end{tabularx}
\textbf{Kontrollfragen:} Gesamtgewinn: 277,30\,\euro{}. Durchschnitt: ca. 46,22\,\euro{}.

\section*{Etappe 7: Diagramm}
Benötigte Bereiche: \code{A3:A8} und \code{G3:G8}. Die höchste Säule sollte bei \textbf{Brezel} liegen, weil die Brezel mit 57,60\,\euro{} den größten Gewinn bringt.

\section*{Etappe 8: Auswertung}
Mögliche Auswertung: Der größte Gesamtgewinn entsteht bei der \textbf{Brezel} mit \textbf{57,60\,\euro{}}. Der höchste Verkaufspreis allein entscheidet nicht, weil auch Einkaufspreis und verkaufte Anzahl wichtig sind. Der Obstbecher kostet zwar 2,20\,\euro{}, bringt insgesamt aber nur 36,30\,\euro{} Gewinn. Für das nächste Schulfest wären besonders Brezeln und Wasser sinnvoll, weil beide hohe Gewinne bringen.

\section*{Etappe 9: Sprinter}
\begin{tabularx}{\textwidth}{|L{5cm}|Y|}
\hline
\rowcolor{gray!20}\textbf{Frage} & \textbf{Antwort}\\\hline
Warum steht \texttt{\$B\$13} in der Formel? & B13 bleibt beim Herunterziehen fest.\\\hline
Welche Formel steht in H3? & \texttt{=C3*(1-\$B\$13)*D3-F3}\\\hline
Gesamtgewinn nach Rabatt & 228,48\,\euro{}\\\hline
Um wie viel Euro sinkt der Gewinn? & 48,82\,\euro{}\\\hline
\end{tabularx}

\begin{tabularx}{\textwidth}{|Y|L{3cm}|}
\hline
\rowcolor{gray!20}\textbf{Artikel} & \textbf{Gewinn nach Rabatt}\\\hline
Wasser & 47,30\,\euro{}\\\hline
Apfelschorle & 37,12\,\euro{}\\\hline
Eistee & 31,32\,\euro{}\\\hline
Brezel & 46,80\,\euro{}\\\hline
Muffin & 36,90\,\euro{}\\\hline
Obstbecher & 29,04\,\euro{}\\\hline
\end{tabularx}
